% GENERATED FILE. Edit canonical lessons, then run npm run generate:books.
% canonical-source-hash: fnv1a64:c49e8a4789956cd1
% canonical-lessons: FR-C529-nappe, FR-C529-couvert, FR-C529-debarrasser, FR-C529-louche, FR-C529-couvercle, FR-R529-revision-522-529

\chapter{À table}
\label{ch:fr-a-table}
\hlchaptermodality{529}

\begin{chapteropening}
\textbf{By the end of this chapter:} \emph{I can talk about laying and clearing up after a meal.}

Complete the last lesson of chapter 529: I can talk about laying and clearing up after a meal.
\end{chapteropening}

\section[la nappe]{la nappe — tablecloth}

\label{lesson:FR-C529-nappe}



\begin{quote}
Before the new one: say the French for sponge, then the French for floor cloth.
\end{quote}

\subsection*{You'll want to know: la nappe}
\textbf{la nappe} — \textquotedblleft{}tablecloth\textquotedblright{}.

\begin{quote}
Pour la fête, ma mère sort la belle nappe blanche.
\end{quote}

\begin{quote}
\emph{For the party my mother gets out the nice white tablecloth.}
\end{quote}

Also a sheet of water or fog: \emph{une nappe de brouillard}.

\begin{cousinweb}
\emph{la nappe} is inherited from Latin \emph{mappa} \textquotedblleft{}a napkin, a cloth\textquotedblright{} (EtymDB; Lewis \& Short), with its first sound changed.
\end{cousinweb}

\begin{grammarlens}[title={What you've built}]
One more word for telling, explaining and describing at B1.
\end{grammarlens}

\subsection*{Your turn}
\begin{itemize}
\raggedright
  \item \emph{Say it:} \emph{la nappe}
  \item \emph{Say it:} \emph{la nappe}, once more
  \item \emph{Say it:} a sentence of your own with \emph{la nappe}
  \item \emph{Recall:} say the French for sponge, then the French for floor cloth, then say \emph{la nappe} again
\end{itemize}

\begin{tcolorbox}[breakable,colback=teal!4,colframe=teal!35!black,title={Before you move on}]
What does la nappe mean? (\textquotedblleft{}Tablecloth\textquotedblright{}.) Say it once more.
\end{tcolorbox}
\section[le couvert]{le couvert — place setting; cutlery}

\label{lesson:FR-C529-couvert}



\begin{quote}
Before the new one: say the French for floor cloth, then the French for tablecloth.
\end{quote}

\subsection*{You'll want to know: le couvert}
\textbf{le couvert} — \textquotedblleft{}place setting; cutlery\textquotedblright{}.

\begin{quote}
Les enfants mettent le couvert avant le dîner.
\end{quote}

\begin{quote}
\emph{The children set out the plates and cutlery before dinner.}
\end{quote}

\emph{Mettre le couvert} — to set out plates and cutlery for a meal. \emph{Les couverts} — the cutlery.

\begin{cousinweb}
\emph{le couvert} is the participle of \emph{couvrir} \textquotedblleft{}to cover\textquotedblright{} (EtymDB), inherited from Latin \emph{cooperīre} \textquotedblleft{}to cover completely\textquotedblright{} (Lewis \& Short).
\end{cousinweb}

\begin{grammarlens}[title={What you've built}]
One more word for telling, explaining and describing at B1.
\end{grammarlens}

\subsection*{Your turn}
\begin{itemize}
\raggedright
  \item \emph{Say it:} \emph{le couvert}
  \item \emph{Say it:} \emph{le couvert}, once more
  \item \emph{Say it:} a sentence of your own with \emph{le couvert}
  \item \emph{Recall:} say the French for floor cloth, then the French for tablecloth, then say \emph{le couvert} again
\end{itemize}

\begin{tcolorbox}[breakable,colback=teal!4,colframe=teal!35!black,title={Before you move on}]
What does le couvert mean? (\textquotedblleft{}Place setting; cutlery\textquotedblright{}.) Say it once more.
\end{tcolorbox}
\section[débarrasser]{débarrasser — to clear away (after a meal); to get rid of}

\label{lesson:FR-C529-debarrasser}



\begin{quote}
Before the new one: say the French for tablecloth, then the French for place setting.
\end{quote}

\subsection*{You'll want to know: débarrasser}
\textbf{débarrasser} — \textquotedblleft{}to clear away (after a meal); to get rid of\textquotedblright{}.

\begin{quote}
Après le repas, mon père débarrasse la table.
\end{quote}

\begin{quote}
\emph{After the meal my father clears the dinner table.}
\end{quote}

\emph{Se débarrasser de} — to get rid of.

\begin{cousinweb}
\emph{débarrasser} is the opposite of \emph{embarrasser} \textquotedblleft{}to hinder\textquotedblright{}. The origin of this root is not settled; a link with \emph{barre} \textquotedblleft{}bar\textquotedblright{} is often proposed.
\end{cousinweb}

\begin{grammarlens}[title={What you've built}]
One more word for telling, explaining and describing at B1.
\end{grammarlens}

\subsection*{Your turn}
\begin{itemize}
\raggedright
  \item \emph{Say it:} \emph{débarrasser}
  \item \emph{Say it:} \emph{débarrasser}, once more
  \item \emph{Say it:} a sentence of your own with \emph{débarrasser}
  \item \emph{Recall:} say the French for tablecloth, then the French for place setting, then say \emph{débarrasser} again
\end{itemize}

\begin{tcolorbox}[breakable,colback=teal!4,colframe=teal!35!black,title={Before you move on}]
What does débarrasser mean? (\textquotedblleft{}To clear away (after a meal); to get rid of\textquotedblright{}.) Say it once more.
\end{tcolorbox}
\section[la louche]{la louche — ladle}

\label{lesson:FR-C529-louche}



\begin{quote}
Before the new one: say the French for place setting, then the French for to clear away.
\end{quote}

\subsection*{You'll want to know: la louche}
\textbf{la louche} — \textquotedblleft{}ladle\textquotedblright{}.

\begin{quote}
Elle sert la soupe avec une louche.
\end{quote}

\begin{quote}
\emph{She serves the soup with a ladle.}
\end{quote}

As an adjective, \emph{louche} means shady: \emph{une affaire louche}. That is a different word.

\begin{cousinweb}
This \emph{louche} \textquotedblleft{}a ladle\textquotedblright{} is usually traced to Frankish \emph{lōtja}, a Germanic word. The adjective \emph{louche} \textquotedblleft{}shady\textquotedblright{} is a different word, from Latin \emph{luscus} \textquotedblleft{}one-eyed\textquotedblright{} (EtymDB).
\end{cousinweb}

\begin{grammarlens}[title={What you've built}]
One more word for telling, explaining and describing at B1.
\end{grammarlens}

\subsection*{Your turn}
\begin{itemize}
\raggedright
  \item \emph{Say it:} \emph{la louche}
  \item \emph{Say it:} \emph{la louche}, once more
  \item \emph{Say it:} a sentence of your own with \emph{la louche}
  \item \emph{Recall:} say the French for place setting, then the French for to clear away, then say \emph{la louche} again
\end{itemize}

\begin{tcolorbox}[breakable,colback=teal!4,colframe=teal!35!black,title={Before you move on}]
What does la louche mean? (\textquotedblleft{}Ladle\textquotedblright{}.) Say it once more.
\end{tcolorbox}
\section[le couvercle]{le couvercle — lid}

\label{lesson:FR-C529-couvercle}



\begin{quote}
Before the new one: say the French for to clear away, then the French for ladle.
\end{quote}

\subsection*{You'll want to know: le couvercle}
\textbf{le couvercle} — \textquotedblleft{}lid\textquotedblright{}.

\begin{quote}
Avec le couvercle sur la casserole, ça cuit plus vite.
\end{quote}

\begin{quote}
\emph{With the lid on the pan, it cooks faster.}
\end{quote}

Built on \emph{couvrir} \textquotedblleft{}to cover\textquotedblright{}.

\begin{cousinweb}
\emph{le couvercle} is inherited from Latin \emph{cooperculum} \textquotedblleft{}a lid\textquotedblright{}, from \emph{cooperīre} \textquotedblleft{}to cover completely\textquotedblright{} (Lewis \& Short).
\end{cousinweb}

\begin{grammarlens}[title={What you've built}]
One more word for telling, explaining and describing at B1.
\end{grammarlens}

\subsection*{Your turn}
\begin{itemize}
\raggedright
  \item \emph{Say it:} \emph{le couvercle}
  \item \emph{Say it:} \emph{le couvercle}, once more
  \item \emph{Say it:} a sentence of your own with \emph{le couvercle}
  \item \emph{Recall:} say the French for to clear away, then the French for ladle, then say \emph{le couvercle} again
\end{itemize}

\begin{tcolorbox}[breakable,colback=teal!4,colframe=teal!35!black,title={Before you move on}]
What does le couvercle mean? (\textquotedblleft{}Lid\textquotedblright{}.) Say it once more.
\end{tcolorbox}
\section[(dialogue)]{Retour sur les chapitres 522–529}

\label{lesson:FR-R529-revision-522-529}



\begin{quote}
Say the French for poetry and comic book.
\end{quote}

\subsection*{Your turn}
Say these aloud:

\begin{itemize}
\raggedright
  \item poetry, comic book, crime novel, biography, tale
  \item episode, season, suspense, (streaming) platform, thrilling
  \item original-language version, to screen, dubbing, (film) director, shooting
  \item to channel-hop, remote control, (television or radio) programme, documentary, to repeat
  \item washing machine, to spin, to hang out, clothes peg, coat hanger
  \item iron, creased, to soak, wardrobe, to hang
  \item to sweep, to dust, to rub, sponge, floor cloth
  \item tablecloth, place setting, to clear away, ladle, lid
\end{itemize}

\begin{tcolorbox}[breakable,colback=teal!4,colframe=teal!35!black,title={Before you move on}]
Poetry? (\textbf{la poésie}.) Comic book? (\textbf{la bande dessinée}.) Crime novel? (\textbf{le polar}.) Biography? (\textbf{la biographie}.) Tale? (\textbf{le conte}.) Episode? (\textbf{l'épisode}.) Season? (\textbf{la saison}.) Suspense? (\textbf{le suspense}.) (streaming) platform? (\textbf{la plateforme}.) Thrilling? (\textbf{palpitant}.) Original-language version? (\textbf{la version originale}.) To screen? (\textbf{projeter}.) Dubbing? (\textbf{le doublage}.) (film) director? (\textbf{le réalisateur}.) Shooting? (\textbf{le tournage}.) To channel-hop? (\textbf{zapper}.) Remote control? (\textbf{la télécommande}.) (television or radio) programme? (\textbf{l'émission}.) Documentary? (\textbf{le documentaire}.) To repeat? (\textbf{rediffuser}.) Washing machine? (\textbf{la machine à laver}.) To spin? (\textbf{essorer}.) To hang out? (\textbf{étendre}.) Clothes peg? (\textbf{la pince à linge}.) Coat hanger? (\textbf{le cintre}.) Iron? (\textbf{le fer à repasser}.) Creased? (\textbf{froissé}.) To soak? (\textbf{tremper}.) Wardrobe? (\textbf{la penderie}.) To hang? (\textbf{accrocher}.) To sweep? (\textbf{balayer}.) To dust? (\textbf{épousseter}.) To rub? (\textbf{frotter}.) Sponge? (\textbf{l'éponge}.) Floor cloth? (\textbf{la serpillière}.) Tablecloth? (\textbf{la nappe}.) Place setting? (\textbf{le couvert}.) To clear away? (\textbf{débarrasser}.) Ladle? (\textbf{la louche}.) Lid? (\textbf{le couvercle}.)
\end{tcolorbox}
